\section{希尔伯特变换}\label{sec:Hilbert}

\subsection{希尔伯特变换}\label{sub:hilbert}

\begin{definition}[希尔伯特变换]
    函数$f:\mathbb{R}\to\mathbb{C}$的\textbf{希尔伯特变换}(Hilbert transform)定义为
    \begin{equation}
        \mathcal{H}f(x)=\frac{1}{\pi}\pv\int_{-\infty}^{\infty}\frac{f(y)}{x-y}\,\dd y=\paren{f*\pv\frac{1}{\pi x}}(x),
    \end{equation}
    称$f$与$\mathcal{H}f$为\textbf{希尔伯特变换对}，记为$f\overset{\mathcal{H}}{\longleftrightarrow}\mathcal{H}f$。
\end{definition}

希尔伯特变换既可以看作积分变换，又可以看作滤波器，因为它是函数与分布$\pv 1/\pi t$的卷积（多数文献中略去了$\pv$的符号，但定义是一致的）。希尔伯特变换在信号处理领域有着重要的应用，例如构造解析信号，将在本小节最后介绍。

\begin{remark}
    一些文献中\cite{brad_osgood_ee261}，希尔伯特变换为$-\pv 1/\pi t$与$f(t)$的卷积，我们很快将看到这就是本书中的希尔伯特逆变换，它们没有本质区别；一些文献中\cite{signal_systems}，直接将希尔伯特变换记为$\hat{f}$，这样做有一定的好处，例如书写简便，但它与傅里叶变换的记号冲突，因此本书中不采用这种记号。
\end{remark}
\begin{remark}
    我们没有对$f$所属的函数空间做出限制。实际上，希尔伯特变换可以定义在多种函数空间上，经典的希尔伯特变换建立在$L^1(\mathbb{R})$上，并且要求函数$f$是赫德尔连续的（我们没有介绍这个概念），希尔伯特变换还可以定义在$L^p(\mathbb{R})\ (1<p<\infty)$上，或者更广义的分布空间上。专著\cite{hilbert_transform}详细讨论了$f$的范围。
\end{remark}

$\displaystyle\pv\frac{1}{x}$作为缓增的正则分布，我们在\refsub{normal_distribution}中推导过其傅里叶变换：
\begin{lralign}
    \mathcal{F}\sbr{\pv\frac{1}{\pi x}}(\xi)=-\im\,\sgn(\xi).
    \switch
    \mathcal{F}\sbr{\pv\frac{1}{\pi x}}(\omega)=-\im\,\sgn(\omega).
\end{lralign}

借助这个结论，我们可以通过傅里叶变换来研究希尔伯特变换的性质。由分布与函数的卷积定理（\cref{thm:distribution_func_conv}），可以得到：
\begin{equation}
    \mathcal{F}\sbr{\mathcal{H}f}(\xi)=-\im\,\sgn(\xi)\mathcal{F}f(\xi).
\end{equation}
可以看到，希尔伯特变换在频域上对应乘以$-\im\,\sgn(\xi)$，这表明希尔伯特变换是一个全通滤波器，它将正频率成分的相位延迟$\pi/2$，将负频率成分的相位超前$\pi/2$，而不改变任意频率成分的幅度。

如果连续取两次希尔伯特变换，在频域上就对应乘以$(-\im\,\sgn(\xi))^2=-1$，也就是：
\begin{equation}
    \mathcal{H}\mathcal{H}f=-f,
\end{equation}
因此，希尔伯特变换的逆变换为：
\begin{definition}[希尔伯特逆变换]
    函数$f:\mathbb{R}\to\mathbb{C}$的\textbf{希尔伯特逆变换}定义为
    \begin{equation}
        \mathcal{H}^{-1}f(x)=-\mathcal{H}f(x)=-\frac{1}{\pi}\pv\int_{-\infty}^{\infty}\frac{f(y)}{x-y}\,\dd y=-\paren{f*\pv\frac{1}{\pi x}}(x).
    \end{equation}
\end{definition}
至此可以看出，希尔伯特变换对是一一的，因此记号$f\overset{\mathcal{H}}{\longleftrightarrow}\mathcal{H}f$是合理的。

下面给出一些希尔伯特变换的计算例子，展示如何使用傅里叶变换来计算希尔伯特变换，以及如何直接使用定义来计算希尔伯特变换。

\begin{example}[三角函数的希尔伯特变换]
    三角函数的希尔伯特变换可以通过它们的傅里叶变换来求解。考虑$\cos(2\pi f_0 x)=\cos(\omega_0 x)$，我们知道：
    \begin{lralign}
        \mathcal{F}\sbr{\cos(2\pi f_0 t)}(\xi)&=\frac{1}{2}(\delta_{f_0}+\delta_{-f_0}),
        \switch
        \mathcal{F}\sbr{\cos(\omega_0 t)}(\xi)&=\pi(\delta_{\omega_0}+\delta_{-\omega_0}),
    \end{lralign}
    因此余弦函数的希尔伯特变换的傅里叶变换为（假设$2\pi f_0=\omega_0>0$）：
    \begin{lralign}
        &\phantom{=}\mathcal{F}\sbr{\mathcal{H}\cos(2\pi f_0 t)}(\xi)\\
        &=-\im\,\sgn(\xi)\cdot\frac{1}{2}(\delta_{f_0}+\delta_{-f_0})\\
        &=-\frac{\im}{2}(\delta_{f_0}-\delta_{-f_0}),
        \switch
        &\phantom{=}\mathcal{F}\sbr{\mathcal{H}\cos(\omega_0 t)}(\omega)\\
        &=-\im\,\sgn(\omega)\cdot\pi(\delta_{\omega_0}+\delta_{-\omega_0})\\
        &=-\im\pi(\delta_{\omega_0}-\delta_{-\omega_0}),
    \end{lralign}
    这正是$\sin(2\pi f_0 t)=\sin(\omega_0 t)$的傅里叶变换，因此我们得到：
    \begin{equation}
        \mathcal{H}\cos(\omega_0 t)=\sin(\omega_0 t),
    \end{equation}
    同理可得：
    \begin{equation}
        \mathcal{H}\sin(\omega_0 t)=-\cos(\omega_0 t).
    \end{equation}
\end{example}

\begin{example}[直接计算希尔伯特变换]
    对于多数函数而言，计算希尔伯特变换时，直接使用定义进行计算比较困难。这里给出一个简单的函数，它能够使用定义计算希尔伯特变换。考虑脉冲宽度为$T$的矩形函数$f(x)=\chi_{[-T/2,T/2]}(x)$，则有：
    \begin{align}
        \mathcal{H}f(x)&=\frac{1}{\pi}\pv\int_{-T/2}^{T/2}\frac{1}{x-y}\,\dd y\nonumber\\
        &=\frac{1}{\pi}\lim_{\epsilon\to 0^+}\paren{\int_{-T/2}^{x-\epsilon}\frac{1}{x-y}\,\dd y+\int_{x+\epsilon}^{T/2}\frac{1}{x-y}\,\dd y}\nonumber\\
        &=\frac{1}{\pi}\lim_{\epsilon\to 0^+}\paren{\ln|x-y|\Big|_{y=-T/2}^{y=x-\epsilon}+\ln|x-y|\Big|_{y=x+\epsilon}^{y=T/2}}\nonumber\\
        &=\frac{1}{\pi}\lim_{\epsilon\to 0^+}\paren{\ln\epsilon-\ln|x+T/2|+\ln|T/2-x|-\ln\epsilon}\nonumber\\
        &=\frac{1}{\pi}\ln\abs{\frac{T/2-x}{x+T/2}}.
    \end{align}
\end{example}

\begin{example}[三角函数的希尔伯特变换，直接计算]
    利用狄利克雷积分（\cref{thm:dirichlet_int}）：
    \begin{equation}
        \int_{-\infty}^{\infty}\frac{\sin ax}{x}\,\dd x=\pi,\quad a>0
    \end{equation}
    也可以直接计算三角函数的希尔伯特变换，这里以正弦函数为例：
    \begin{align}
        \mathcal{H}\sin(\omega_0 t)&=\frac{1}{\pi}\pv\int_{-\infty}^{\infty}\frac{\sin(\omega_0 y)}{t-y}\,\dd y\nonumber\\
        &=\frac{1}{\pi}\lim_{\epsilon\to 0^+}\paren{\int_{-\infty}^{t-\epsilon}\frac{\sin(\omega_0 y)}{t-y}\,\dd y+\int_{t+\epsilon}^{\infty}\frac{\sin(\omega_0 y)}{t-y}\,\dd y}\nonumber\\
        &=\frac{1}{\pi}\lim_{\epsilon\to 0^+}\paren{\int_{\epsilon}^{\infty}\frac{\sin(\omega_0(t-u))}{u}\,\dd u+\int_{-\infty}^{-\epsilon}\frac{\sin(\omega_0(t-u))}{u}\,\dd u}&(u=t-y)\nonumber\\
        &=\frac{1}{\pi}\lim_{\epsilon\to 0^+}\left(\int_{\epsilon}^{\infty}\frac{\sin(\omega_0 t)\cos(\omega_0 u)-\cos(\omega_0 t)\sin(\omega_0 u)}{u}\,\dd u\right.\nonumber\\
        &\qquad\left.{}+\int_{-\infty}^{-\epsilon}\frac{\sin(\omega_0 t)\cos(\omega_0 u)-\cos(\omega_0 t)\sin(\omega_0 u)}{u}\,\dd u\right)\nonumber\\
        &=\frac{1}{\pi}\lim_{\epsilon\to 0^+}\left(\sin(\omega_0 t)\int_{\epsilon}^{\infty}\frac{\cos(\omega_0 u)}{u}\,\dd u+\sin(\omega_0 t)\int_{-\infty}^{-\epsilon}\frac{\cos(\omega_0 u)}{u}\,\dd u\right.\nonumber\\
        &\qquad\left.{}-\cos(\omega_0 t)\int_{\epsilon}^{\infty}\frac{\sin(\omega_0 u)}{u}\,\dd u-\cos(\omega_0 t)\int_{-\infty}^{-\epsilon}\frac{\sin(\omega_0 u)}{u}\,\dd u\right)\nonumber\\
        &=\frac{1}{\pi}\lim_{\epsilon\to 0^+}\paren{-\cos(\omega_0 t)\int_{\epsilon}^{\infty}\frac{\sin(\omega_0 u)}{u}\,\dd u+\cos(\omega_0 t)\int_{\epsilon}^{\infty}\frac{\sin(\omega_0 u)}{u}\,\dd u}\nonumber\\
        &=-\cos(\omega_0 t).
    \end{align}
    倒数第三行到倒数第二行运用了函数的奇偶性。
\end{example}

\subsection{希尔伯特变换的性质}\label{sub:prop_of_hilbert}

简洁起见，本小节主要使用频率形式的傅里叶变换。如果$f\in L^2(\mathbb{R})$，即$f(t)$是一个能量有限信号，则$\mathcal{H}f\in L^2(\mathbb{R})$，并且$\norm{\mathcal{H}f}_2=\norm{f}_2$，进一步，$L^2(\mathbb{R})$上的\textbf{希尔伯特变换是保内积的}：
\begin{theorem}
    设$f,g\in L^2(\mathbb{R})$，则有
    \begin{equation}
        (\mathcal{H}f,\mathcal{H}g)=(f,g),
    \end{equation}
    特别地，$\norm{\mathcal{H}f}_2=\norm{f}_2$，这称之为希尔伯特变换的\textbf{等能量性}。
\end{theorem}
\begin{proof}
    由希尔伯特变换的傅里叶变换表达式，有
    \begin{align}
        (\mathcal{H}f,\mathcal{H}g)&=\int_{-\infty}^{\infty}\mathcal{H}f(t)\overline{\mathcal{H}g(t)}\,\dd t\nonumber\\
        &=\int_{-\infty}^{\infty}\mathcal{F}\sbr{\mathcal{H}f}(\xi)\overline{\mathcal{F}\sbr{\mathcal{H}g}(\xi)}\,\dd\xi\nonumber\\
        &=\int_{-\infty}^{\infty}\paren{-\im\,\sgn(\xi)\mathcal{F}f(\xi)}\overline{-\im\,\sgn(\xi)\mathcal{F}g(\xi)}\,\dd\xi\nonumber\\
        &=\int_{-\infty}^{\infty}\mathcal{F}f(\xi)\overline{\mathcal{F}g(\xi)}\,\dd\xi\nonumber\\
        &=(f,g).
    \end{align}
\end{proof}

希尔伯特变换还具有\textbf{正交性}：
\begin{theorem}
    设实函数$f\in L^2(\mathbb{R})$，则有
    \begin{equation}
        (f,\mathcal{H}f)=0.
    \end{equation}
\end{theorem}
由于这个原因，有时将$\mathcal{H}f$称为$f$的\textbf{正交伴随}(orthogonal complement)。
\begin{proof}
    由希尔伯特变换的傅里叶变换表达式，有
    \begin{align}
        (f,\mathcal{H}f)&=\int_{-\infty}^{\infty}f(t)\overline{\mathcal{H}f(t)}\,\dd t\nonumber\\
        &=\int_{-\infty}^{\infty}\mathcal{F}f(\xi)\overline{\mathcal{F}\sbr{\mathcal{H}f}(\xi)}\,\dd\xi\nonumber\\
        &=\int_{-\infty}^{\infty}\mathcal{F}f(\xi)\overline{-\im\,\sgn(\xi)\mathcal{F}f(\xi)}\,\dd\xi\nonumber\\
        &=\im\int_{-\infty}^{\infty}\sgn(\xi)\abs{\mathcal{F}f(\xi)}^2\,\dd\xi\nonumber\\
        &=0.
    \end{align}
    最后一个等号是因为实函数的傅里叶变换具有共轭对称性，$\abs{\mathcal{F}f(\xi)}$是偶函数，因此$\sgn(\xi)\abs{\mathcal{F}f(\xi)}^2$是奇函数。
\end{proof}

此外，希尔伯特变换具有奇偶对称性：
\begin{theorem}
    $f$与$\mathcal{H}f$的奇偶性相反，即：
    \begin{itemize}
        \item 如果$f$是偶函数，则$\mathcal{H}f$是奇函数；
        \item 如果$f$是奇函数，则$\mathcal{H}f$是偶函数。
    \end{itemize}
\end{theorem}
\begin{proof}
    根据傅里叶变换的对称性，如果$f$是偶函数，则$\mathcal{F}f$也是偶函数，于是
    \begin{align}
        \mathcal{F}\sbr{\mathcal{H}f}(\xi)&=-\im\,\sgn(\xi)\mathcal{F}f(\xi),\\
        \mathcal{F}\sbr{\mathcal{H}f}(-\xi)&=-\im\,\sgn(-\xi)\mathcal{F}f(-\xi)\nonumber\\
        &=\im\,\sgn(\xi)\mathcal{F}f(\xi)\nonumber\\
        &=-\mathcal{F}\sbr{\mathcal{H}f}(\xi),
    \end{align}
    即$\mathcal{F}\sbr{\mathcal{H}f}(\xi)$是奇函数，因此$\mathcal{H}f$是奇函数。

    类似地，如果$f$是奇函数，则$\mathcal{F}f$也是奇函数，$\mathcal{F}\sbr{\mathcal{H}f}(\xi)=-\im\,\sgn(\xi)\mathcal{F}f(\xi)$是偶函数，因此$\mathcal{H}f$是偶函数。直接使用希尔伯特变换的定义也可验证这一性质。
\end{proof}

\subsection{解析信号}\label{sub:analytic_signal}

希尔伯特变换可以用来构造解析信号。解析信号在通信领域有着重要的应用，例如调制解调、单边带调制等。
\begin{definition}[解析信号]
    函数$f\in L^1(\mathbb{R})$称为\textbf{解析信号}(analytic signal)，如果它仅含正频率成分，即
    \begin{equation}
        \supp(\mathcal{F}f)\subseteq[0,\infty).
    \end{equation}
\end{definition}

实函数的傅里叶变换满足共轭对称性：$\mathcal{F}f^-=\overline{\mathcal{F}f}$，即幅度谱为偶函数而相位谱为奇函数，因此实函数不可能是解析信号。但如果考虑$2u\cdot\mathcal{F}f$，其中$u$为单位阶跃函数，则该信号仅含正频率成分（我们马上将看到为什么是两倍的$u\mathcal{F}f$）。

对$2u\cdot\mathcal{F}f$做傅里叶逆变换，并记相应的时域函数为$\mathcal{Z}_f(t)$，则有：
\begin{align}
    \mathcal{Z}_f(t)&=\mathcal{F}^{-1}\sbr{2u(\xi)\cdot\mathcal{F}f(\xi)}(t)\nonumber\\
    &=\mathcal{F}^{-1}\sbr{(1+\sgn(\xi))\mathcal{F}f(\xi)}(t)\nonumber\\
    &=\mathcal{F}\mathcal{F}^{-1}f(t)+\mathcal{F}^{-1}\sbr{\sgn(\xi)\mathcal{F}f(\xi)}(t)\nonumber\\
    &=f(t)+\im\mathcal{H}f(t).
\end{align}

显然，实函数的希尔伯特变换仍是实函数。因此可以说，\textbf{在实函数的虚部中加入它的希尔伯特变换，就得到了一个解析信号}，原有的实函数$f(t)$是其实部。当然，以上解析信号的构造过程对于复值函数同样有效，只是复值函数无法通过对解析信号取实部来恢复。

\begin{definition}
    设$f\in L^1(\mathbb{R})$，定义
    \begin{equation}
        \mathcal{Z}_f(t)=f(t)+\im\mathcal{H}f(t)
    \end{equation}
    称之为\textbf{与}$f$\textbf{相关的解析信号}。
    频域上，$\mathcal{Z}_f$是由$f$截去负频率成分，而正频率成分加倍得到的。
\end{definition}

与解析信号相关的解析信号并不是它本身，而是它的两倍，这一点从频域上是容易看出来的。

对于一个解析信号$g(t)$，它的实部为：
\begin{equation}
    \Re(g(t))=\frac{g(t)+\overline{g(t)}}{2}.
\end{equation}
实部的傅里叶变换为：
\begin{equation}
    \mathcal{F}\sbr{\Re(g(t))}=\frac{\mathcal{F}g(\xi)+\mathcal{F}\overline{g}(\xi)}{2}=\frac{\mathcal{F}g(\xi)+\mathcal{F}g(-\xi)}{2}.
\end{equation}
可见取解析信号的实部相当于在频域上做偶延拓再将幅值减半，这正是由实函数构造与之相关的解析信号的逆过程。因此，\textbf{一个信号是解析信号，当且仅当它的虚部是实部的希尔伯特变换}。

以上构造表明：若某线性时不变系统的响应恰好是激励的希尔伯特变换，就可借此构造与激励相关的解析信号。然而，希尔伯特变换的单位冲激响应为$1/\pi t$，其本身非因果，因此对应系统不可实现。

\subsection{希尔伯特变换关系}\label{sub:hilbert_relation}

解析信号在频域中仅含正频率成分，可视为“频域因果”信号；相应地，其时域实部与虚部构成希尔伯特变换对。由此可类比猜想：若线性时不变系统在时域因果，则其频率响应的实部与虚部也应满足同样的希尔伯特变换关系。下面给出证明。

设$h(t)$是因果线性时不变系统的单位冲激响应，频率响应为$H(\omega)$。由$h(t)=u(t)h(t)$，可得：
\begin{align}
    H(\omega)&=\mathcal{F}\sbr{h(t)u(t)}(\omega)\nonumber\\
    &=\frac{1}{2\pi}H(\omega)*\paren{\pi\delta(\omega)+\pv\frac{1}{\im\omega}}\nonumber\\
    &=\frac{1}{2}H(\omega)+\pv\frac{1}{2\pi\im\omega}*H(\omega).
\end{align}

考察上式的实部和虚部，记$H(\omega)=R(\omega)+\im I(\omega)$，则有：
\begin{align}
    R(\omega)&=\frac{1}{2}R(\omega)-\pv\frac{1}{2\pi\omega}*I(\omega),\\
    I(\omega)&=\frac{1}{2}I(\omega)+\pv\frac{1}{2\pi\omega}*R(\omega).
\end{align}
整理以上两式即得：
\begin{theorem}[希尔伯特变换关系]
    设$h(t)$为因果线性时不变系统的单位冲激响应，$H(\omega)=R(\omega)+\im I(\omega)$为系统的频率响应，则$R(\omega)$与$I(\omega)$构成希尔伯特变换对，即：
    \begin{align}
        R(\omega)&=-\pv\frac{1}{\pi\omega}*I(\omega)=\mathcal{H}^{-1}I(\omega),\\
        I(\omega)&=\pv\frac{1}{\pi\omega}*R(\omega)=\mathcal{H}R(\omega).
    \end{align}
    将这种关系称为\textbf{希尔伯特变换关系}(Hilbert transform relation)。
\end{theorem}
